\documentclass[11pt,letterpaper]{article}
\usepackage[utf8]{inputenc}
\usepackage[T1]{fontenc}
\usepackage[spanish,es-noshorthands,es-tabla]{babel}
\usepackage{lmodern}
\usepackage[margin=2.2cm]{geometry}
\usepackage{amsmath,amssymb,amsthm,mathtools}
\usepackage{enumitem}
\usepackage{booktabs,array,tabularx}
\usepackage{multicol}
\usepackage{xcolor}
\usepackage[most]{tcolorbox}
\usepackage{tikz}
\usetikzlibrary{shapes.geometric,arrows.meta,positioning,calc}
\usepackage{fancyhdr}
\usepackage{titlesec}
\usepackage{listings}
\usepackage[hidelinks]{hyperref}

\definecolor{azul}{HTML}{1F4E79}
\definecolor{azulclaro}{HTML}{E8F0F8}
\definecolor{verde}{HTML}{2E7D32}
\definecolor{verdeclaro}{HTML}{EAF5EA}
\definecolor{naranja}{HTML}{B85C00}
\definecolor{naranjaclaro}{HTML}{FFF3E0}
\definecolor{gris}{HTML}{555555}
\definecolor{rojo}{HTML}{B71C1C}

\titleformat{\section}{\Large\bfseries\color{azul}}{\thesection}{0.6em}{}
\titleformat{\subsection}{\large\bfseries\color{azul}}{\thesubsection}{0.5em}{}
\titleformat{\subsubsection}{\normalsize\bfseries\color{gris}}{\thesubsubsection}{0.5em}{}

\setlength{\parindent}{0pt}
\setlength{\parskip}{4pt}
\setlist{itemsep=2pt,topsep=3pt}

\pagestyle{fancy}
\fancyhf{}
\renewcommand{\headrulewidth}{0.4pt}

% ---------- Entornos ----------
\newcounter{ej}
\newtcolorbox{ejercicio}[2][]{enhanced,breakable,colback=white,colframe=azul,
  boxrule=0.6pt,arc=2pt,left=6pt,right=6pt,top=4pt,bottom=4pt,
  title={\textbf{#2}},fonttitle=\small\bfseries,coltitle=white,colbacktitle=azul,#1}
\newtcolorbox{solucion}[2][]{enhanced,breakable,colback=verdeclaro,colframe=verde,
  boxrule=0.6pt,arc=2pt,left=6pt,right=6pt,top=4pt,bottom=4pt,
  title={\textbf{#2}},fonttitle=\small\bfseries,coltitle=white,colbacktitle=verde,#1}
\newtcolorbox{nota}[1][]{enhanced,breakable,colback=naranjaclaro,colframe=naranja,
  boxrule=0.5pt,arc=2pt,left=6pt,right=6pt,top=3pt,bottom=3pt,#1}
\newtcolorbox{formulario}[1][]{enhanced,breakable,colback=azulclaro,colframe=azul,
  boxrule=0.5pt,arc=2pt,left=6pt,right=6pt,top=3pt,bottom=3pt,#1}

\newcommand{\dif}[1]{%
  \ifcase#1\or\textcolor[HTML]{66BB6A}{$\bigstar$}\or\textcolor[HTML]{FFA726}{$\bigstar\bigstar$}\or\textcolor[HTML]{EF5350}{$\bigstar\bigstar\bigstar$}\fi}
\newcommand{\fuente}[1]{\hfill{\footnotesize\textcolor{white!80!azul}{[#1]}}}
\newcommand{\ptst}[1]{{\footnotesize\textcolor{white!80!azul}{(#1 pts)}}}
\newcommand{\pts}[1]{\textcolor{gris}{(#1 pts)}}

% ---------- Lógica ----------
\newcommand{\imp}{\Rightarrow}
\newcommand{\sii}{\Leftrightarrow}
\newcommand{\ent}{\models}
\newcommand{\der}{\vdash}
\newcommand{\p}[1]{\ensuremath{\textit{#1}}}

% ---------- Árboles de juego ----------
\tikzset{
  maxn/.style={draw,regular polygon,regular polygon sides=3,minimum size=0.75cm,inner sep=0pt,fill=azulclaro},
  minn/.style={draw,regular polygon,regular polygon sides=3,shape border rotate=180,minimum size=0.75cm,inner sep=0pt,fill=naranjaclaro},
  chn/.style={draw,circle,minimum size=0.6cm,inner sep=0pt,fill=verdeclaro},
  hoja/.style={draw,rectangle,minimum size=0.5cm,inner sep=1pt,font=\small},
  gnode/.style={draw,circle,minimum size=0.75cm,inner sep=0pt,font=\small\bfseries},
  garc/.style={-{Stealth[length=2mm]},thick},
}

\lstset{basicstyle=\ttfamily\small,columns=fullflexible,frame=single,framerule=0.3pt,rulecolor=\color{gris},
  literate={á}{{\'a}}1 {é}{{\'e}}1 {í}{{\'i}}1 {ó}{{\'o}}1 {ú}{{\'u}}1 {ñ}{{\~n}}1}

% ================= Figuras compartidas (guía y respuestas) =================

% ---- Gráfica G1: búsqueda no informada / informada ----
\newcommand{\figGUno}{%
\begin{tikzpicture}[scale=1.0]
  \node[gnode] (S) at (0,0) {S};
  \node[gnode] (A) at (2.4,1.3) {A};
  \node[gnode] (B) at (2.4,-1.3) {B};
  \node[gnode] (C) at (4.9,1.3) {C};
  \node[gnode] (D) at (4.9,-1.3) {D};
  \node[gnode,fill=verdeclaro] (G) at (7.3,0) {G};
  \draw[garc] (S) -- node[above left]{2} (A);
  \draw[garc] (S) -- node[below left]{1} (B);
  \draw[garc] (A) -- node[above]{2} (C);
  \draw[garc] (A) -- node[pos=0.3,right]{4} (D);
  \draw[garc] (B) -- node[below]{1} (D);
  \draw[garc] (B) to[bend right=50] node[below]{9} (G);
  \draw[garc] (C) -- node[above right]{4} (G);
  \draw[garc] (D) -- node[below right]{3} (G);
\end{tikzpicture}\\[2pt]
{\small\begin{tabular}{c|cccccc}$n$ & S & A & B & C & D & G\\\hline $h(n)$ & 4 & 2 & 3 & 1 & 2 & 0\end{tabular}}}

% ---- Gráfica G2: heurística admisible pero inconsistente ----
\newcommand{\figGDos}{%
\begin{tikzpicture}
  \node[gnode] (S) at (0,0) {S};
  \node[gnode] (A) at (2,1) {A};
  \node[gnode] (B) at (2,-1) {B};
  \node[gnode] (C) at (4,0) {C};
  \node[gnode,fill=verdeclaro] (G) at (6,0) {G};
  \draw[garc] (S) -- node[above left]{1} (A);
  \draw[garc] (S) -- node[below left]{1} (B);
  \draw[garc] (A) -- node[above right]{1} (C);
  \draw[garc] (B) -- node[below right]{2} (C);
  \draw[garc] (C) -- node[above]{3} (G);
\end{tikzpicture}\\[2pt]
{\small\begin{tabular}{c|ccccc}$n$ & S & A & B & C & G\\\hline $h(n)$ & 2 & 4 & 1 & 1 & 0\end{tabular}}}

% ---- Gráfica G3: examen simulado ----
\newcommand{\figGTres}{%
\begin{tikzpicture}
  \node[gnode] (S) at (0,0) {S};
  \node[gnode] (A) at (2.2,1.2) {A};
  \node[gnode] (B) at (2.2,-1.2) {B};
  \node[gnode] (C) at (4.6,0) {C};
  \node[gnode,fill=verdeclaro] (G) at (6.8,0) {G};
  \draw[garc] (S) -- node[above left]{1} (A);
  \draw[garc] (S) -- node[below left]{4} (B);
  \draw[garc] (A) -- node[right]{2} (B);
  \draw[garc] (A) -- node[above right]{5} (C);
  \draw[garc] (B) -- node[below right]{1} (C);
  \draw[garc] (C) -- node[above]{3} (G);
  \draw[garc] (B) to[bend right=35] node[below]{6} (G);
\end{tikzpicture}\\[2pt]
{\small\begin{tabular}{c|ccccc}$n$ & S & A & B & C & G\\\hline $h(n)$ & 6 & 5 & 2 & 3 & 0\end{tabular}}}

% ---- Gráfica bidireccional ----
\newcommand{\figBidir}{%
\begin{tikzpicture}
  \node[gnode] (S) at (0,0) {S};
  \node[gnode] (M) at (3.5,1.3) {M};
  \node[gnode] (A) at (2.2,-1.2) {A};
  \node[gnode] (B) at (4.8,-1.2) {B};
  \node[gnode,fill=verdeclaro] (G) at (7,0) {G};
  \draw[thick] (S) -- node[above left]{10} (M);
  \draw[thick] (M) -- node[above right]{10} (G);
  \draw[thick] (S) -- node[below left]{8} (A);
  \draw[thick] (A) -- node[below]{3} (B);
  \draw[thick] (B) -- node[below right]{8} (G);
\end{tikzpicture}}

% ---- Árbol T1 (MAX-MIN-MAX-hojas) ----
\newcommand{\figTUno}{%
\begin{tikzpicture}[level distance=1.25cm,
  level 1/.style={sibling distance=4.6cm},
  level 2/.style={sibling distance=2.3cm},
  level 3/.style={sibling distance=1.05cm}]
  \node[maxn] {R}
    child {node[minn] {A}
      child {node[maxn] {\scriptsize A1} child {node[hoja]{4}} child {node[hoja]{7}}}
      child {node[maxn] {\scriptsize A2} child {node[hoja]{2}} child {node[hoja]{9}}}}
    child {node[minn] {B}
      child {node[maxn] {\scriptsize B1} child {node[hoja]{6}} child {node[hoja]{1}}}
      child {node[maxn] {\scriptsize B2} child {node[hoja]{3}} child {node[hoja]{8}}}}
    child {node[minn] {C}
      child {node[maxn] {\scriptsize C1} child {node[hoja]{5}} child {node[hoja]{10}}}
      child {node[maxn] {\scriptsize C2} child {node[hoja]{12}} child {node[hoja]{2}}}};
\end{tikzpicture}}

% ---- Árbol T2 (AIMA Fig. 6.2) ----
\newcommand{\figTDos}{%
\begin{tikzpicture}[level distance=1.3cm,
  level 1/.style={sibling distance=3.6cm},
  level 2/.style={sibling distance=1.0cm}]
  \node[maxn] {R}
    child {node[minn] {A} child {node[hoja]{3}} child {node[hoja]{12}} child {node[hoja]{8}}}
    child {node[minn] {B} child {node[hoja]{2}} child {node[hoja]{4}} child {node[hoja]{6}}}
    child {node[minn] {C} child {node[hoja]{14}} child {node[hoja]{5}} child {node[hoja]{2}}};
\end{tikzpicture}}

% ---- Árbol T3 (profundidad 4) ----
\newcommand{\figTTres}{%
\begin{tikzpicture}[level distance=1.2cm,
  level 1/.style={sibling distance=7.6cm},
  level 2/.style={sibling distance=3.8cm},
  level 3/.style={sibling distance=1.9cm},
  level 4/.style={sibling distance=0.9cm}]
  \node[maxn] {R}
    child {node[minn] {\scriptsize 1}
      child {node[maxn] {\scriptsize 11}
        child {node[minn] {\tiny 111} child {node[hoja]{2}} child {node[hoja]{5}}}
        child {node[minn] {\tiny 112} child {node[hoja]{8}} child {node[hoja]{1}}}}
      child {node[maxn] {\scriptsize 12}
        child {node[minn] {\tiny 121} child {node[hoja]{4}} child {node[hoja]{6}}}
        child {node[minn] {\tiny 122} child {node[hoja]{7}} child {node[hoja]{3}}}}}
    child {node[minn] {\scriptsize 2}
      child {node[maxn] {\scriptsize 21}
        child {node[minn] {\tiny 211} child {node[hoja]{9}} child {node[hoja]{2}}}
        child {node[minn] {\tiny 212} child {node[hoja]{1}} child {node[hoja]{6}}}}
      child {node[maxn] {\scriptsize 22}
        child {node[minn] {\tiny 221} child {node[hoja]{5}} child {node[hoja]{8}}}
        child {node[minn] {\tiny 222} child {node[hoja]{3}} child {node[hoja]{4}}}}};
\end{tikzpicture}}

% ---- Árbol expectiminimax ----
\newcommand{\figExp}{%
\begin{tikzpicture}[level distance=1.3cm,
  level 1/.style={sibling distance=4.6cm},
  level 2/.style={sibling distance=2.3cm},
  level 3/.style={sibling distance=1.0cm}]
  \node[maxn] {R}
    child {node[chn] {$c_1$}
      child {node[minn]{} child {node[hoja]{3}} child {node[hoja]{9}} edge from parent node[left,font=\scriptsize]{0.5}}
      child {node[minn]{} child {node[hoja]{6}} child {node[hoja]{4}} edge from parent node[right,font=\scriptsize]{0.5}}}
    child {node[chn] {$c_2$}
      child {node[minn]{} child {node[hoja]{2}} child {node[hoja]{10}} edge from parent node[left,font=\scriptsize]{0.5}}
      child {node[minn]{} child {node[hoja]{20}} child {node[hoja]{15}} edge from parent node[right,font=\scriptsize]{0.5}}}
    child {node[chn] {$c_3$}
      child {node[minn]{} child {node[hoja]{7}} child {node[hoja]{8}} edge from parent node[left,font=\scriptsize]{0.3}}
      child {node[minn]{} child {node[hoja]{4}} child {node[hoja]{5}} edge from parent node[right,font=\scriptsize]{0.7}}};
\end{tikzpicture}}

% ---- Árbol de tres jugadores ----
\newcommand{\figMaxN}{%
\begin{tikzpicture}[level distance=1.3cm,
  level 1/.style={sibling distance=6.2cm},
  level 2/.style={sibling distance=3.1cm},
  level 3/.style={sibling distance=1.55cm},
  every node/.style={font=\small}]
  \node[draw,circle,fill=azulclaro] {1}
    child {node[draw,circle,fill=naranjaclaro] {2}
      child {node[draw,circle,fill=verdeclaro] {3} child {node{$(1,2,6)$}} child {node{$(4,2,3)$}}}
      child {node[draw,circle,fill=verdeclaro] {3} child {node{$(6,1,2)$}} child {node{$(7,4,1)$}}}}
    child {node[draw,circle,fill=naranjaclaro] {2}
      child {node[draw,circle,fill=verdeclaro] {3} child {node{$(5,1,1)$}} child {node{$(3,5,2)$}}}
      child {node[draw,circle,fill=verdeclaro] {3} child {node{$(7,7,1)$}} child {node{$(5,4,5)$}}}};
\end{tikzpicture}}

% ---- Paisaje 1D ----
\newcommand{\figPaisaje}{%
\begin{tikzpicture}[yscale=0.35,xscale=0.62]
  \foreach \x/\y in {0/3,1/5,2/6,3/5,4/4,5/6,6/8,7/9,8/7,9/4,10/2,11/5,12/7,13/6,14/3}{
    \fill[azul!60] (\x-0.35,0) rectangle (\x+0.35,\y);
    \node[font=\scriptsize] at (\x,\y+0.6) {\y};
    \node[font=\scriptsize,text=gris] at (\x,-0.7) {\x};
  }
  \draw (-0.6,0) -- (14.6,0);
  \node[font=\scriptsize,anchor=east] at (-0.7,-0.7) {estado};
  \node[font=\scriptsize,anchor=east] at (-0.7,4) {$f$};
\end{tikzpicture}}

% ---- Tablero 4 reinas [2,1,4,4] ----
\newcommand{\figReinas}[4]{%
\begin{tikzpicture}[scale=0.55]
  \foreach \i in {0,...,3}{\foreach \j in {0,...,3}{
    \pgfmathparse{mod(\i+\j,2)==0 ? "gris!25" : "white"}
    \fill[\pgfmathresult] (\i,\j) rectangle (\i+1,\j+1);}}
  \draw (0,0) rectangle (4,4);
  \foreach \c/\r in {0/#1,1/#2,2/#3,3/#4}{\node at (\c+0.5,\r-0.5) {\Large$\mathbf{Q}$};}
  \foreach \c in {1,...,4}{\node[font=\scriptsize] at (\c-0.5,-0.4) {\c};}
  \foreach \r in {1,...,4}{\node[font=\scriptsize] at (-0.4,\r-0.5) {\r};}
\end{tikzpicture}}

% ---- 8-puzzle ----
\newcommand{\puzzle}[9]{%
\begin{tikzpicture}[scale=0.5]
  \draw (0,0) grid (3,3);
  \node at (0.5,2.5){#1};\node at (1.5,2.5){#2};\node at (2.5,2.5){#3};
  \node at (0.5,1.5){#4};\node at (1.5,1.5){#5};\node at (2.5,1.5){#6};
  \node at (0.5,0.5){#7};\node at (1.5,0.5){#8};\node at (2.5,0.5){#9};
\end{tikzpicture}}

% ---- Gato ----
\newcommand{\gato}[9]{%
\begin{tikzpicture}[scale=0.45]
  \draw (1,0)--(1,3) (2,0)--(2,3) (0,1)--(3,1) (0,2)--(3,2);
  \node at (0.5,2.5){#1};\node at (1.5,2.5){#2};\node at (2.5,2.5){#3};
  \node at (0.5,1.5){#4};\node at (1.5,1.5){#5};\node at (2.5,1.5){#6};
  \node at (0.5,0.5){#7};\node at (1.5,0.5){#8};\node at (2.5,0.5){#9};
\end{tikzpicture}}

% ---- Mapa de Australia (grafo de restricciones) ----
\newcommand{\figAustralia}{%
\begin{tikzpicture}[scale=0.9]
  \node[gnode] (WA) at (0,0.5) {\scriptsize WA};
  \node[gnode] (NT) at (1.6,1.6) {\scriptsize NT};
  \node[gnode] (SA) at (2.0,0) {\scriptsize SA};
  \node[gnode] (Q) at (3.6,1.4) {\scriptsize Q};
  \node[gnode] (NSW) at (4.0,-0.2) {\scriptsize NSW};
  \node[gnode] (V) at (3.0,-1.5) {\scriptsize V};
  \node[gnode] (T) at (4.6,-2.2) {\scriptsize T};
  \draw[thick] (WA)--(NT) (WA)--(SA) (NT)--(SA) (NT)--(Q) (SA)--(Q) (SA)--(NSW) (SA)--(V) (Q)--(NSW) (NSW)--(V);
\end{tikzpicture}}

\fancyhead[L]{\small\textcolor{gris}{Inteligencia Artificial 2027-I · Guía para el Examen 1}}
\fancyhead[R]{\small\textcolor{gris}{Enunciados}}
\fancyfoot[C]{\small\thepage}

\begin{document}

% ======================= PORTADA =======================
\begin{titlepage}
\centering
\vspace*{1.5cm}
{\color{azul}\rule{\linewidth}{1.2pt}}\\[0.6cm]
{\Huge\bfseries\color{azul} Guía de estudio para el Examen 1}\\[0.3cm]
{\Large Ejercicios teóricos de Inteligencia Artificial}\\[0.2cm]
{\large Posgrado en Ciencias e Ingeniería de la Computación, UNAM · 2027-I}\\[0.4cm]
{\color{azul}\rule{\linewidth}{1.2pt}}\\[1.2cm]

\begin{minipage}{0.86\linewidth}
\large
\textbf{Cobertura:} agentes racionales; formulación de problemas; búsqueda no informada
e informada; heurísticas (admisibilidad, consistencia, dominancia); búsqueda
bidireccional; CSP (backtracking, revisión hacia adelante, AC-3, MRV/LCV);
búsqueda local y metaheurísticas; búsqueda con adversarios (minimax, alfa-beta,
funciones de evaluación, expectimax, juegos de $n$ jugadores, utilidad);
agentes lógicos y lógica proposicional (modelos, vinculación, CNF, resolución,
Horn, encadenamiento, DPLL); lógica de primer orden (traducción, CNF,
skolemización, unificación, resolución, Prolog).

\vspace{0.6cm}
\textbf{Fuentes:} tareas 1, 2 y 3 del curso; diapositivas 2–12; notas de clase;
Russell \& Norvig, \emph{Artificial Intelligence: A Modern Approach}, 4.ª ed.\ global
(citado como \textbf{AIMA}), capítulos 2–9 y 15.
\end{minipage}

\vfill
\begin{tabular}{ll}
\dif{1} & Nivel 1 — conceptos y definiciones (respuesta corta)\\[2pt]
\dif{2} & Nivel 2 — ejercicios de procedimiento (trazas, cálculos, conversiones)\\[2pt]
\dif{3} & Nivel 3 — tipo examen: formulaciones, demostraciones y problemas largos\\
\end{tabular}

\vspace{1cm}
{\small Las respuestas completas están en el documento complementario
\emph{Guía para el Examen 1 — Respuestas}. Intenta cada ejercicio antes de consultarlas.}
\end{titlepage}

\tableofcontents
\newpage

% ======================= CÓMO USAR =======================
\section*{Cómo usar esta guía}
\addcontentsline{toc}{section}{Cómo usar esta guía}

\begin{itemize}
  \item Los ejercicios están ordenados \textbf{de menor a mayor dificultad}: Parte A (conceptos),
  Parte B (procedimientos), Parte C (tipo examen) y Parte D (examen simulado de 2 horas).
  Dentro de cada parte se sigue el orden del curso.
  \item El profesor indicó que el examen se parecerá a las \textbf{preguntas teóricas de las
  tareas}. Por ello la Parte C reproduce ese estilo: formular un espacio de estados y acotar su
  tamaño, demostrar admisibilidad/consistencia, formular CSP, analizar alfa-beta formalmente,
  clasificar oraciones, resolver con resolución, traducir a LPO y unificar.
  \item Los problemas marcados con \textbf{[AIMA …]} provienen del libro (ejemplos del texto o
  ejercicios del sitio de ejercicios del libro); los marcados con \textbf{[Clase]} o
  \textbf{[Tarea]} son variantes del material del curso.
  \item Convenciones, salvo que se diga otra cosa: búsqueda en \emph{grafo} (conjunto de
  explorados), \textbf{prueba de meta al extraer el nodo} de la frontera, empates resueltos por
  \textbf{orden alfabético}; en alfa-beta se poda cuando $v\le\alpha$ (en MIN) o $v\ge\beta$ (en MAX).
  \item Un ejercicio de Nivel 3 típico toma entre 10 y 20 minutos. Cronométrate.
  \item \textbf{Parte E (agregada tras ver el examen de 2025):} reproduce el formato real. Si tienes poco
  tiempo, prioriza E1–E5, la Parte D y los ejercicios B15–B19, B25, B29–B30, C3 y C13.
\end{itemize}

\subsection*{Plan de estudio (viernes 25 de septiembre → jueves 1 de octubre)}
\begin{center}
\renewcommand{\arraystretch}{1.25}
\begin{tabularx}{\linewidth}{>{\bfseries}l X r}
\toprule
Día & Actividades & Tiempo \\
\midrule
Vie 25 & Leer el formulario (sección 1). Parte A: A1–A3. Parte B: B1–B7 (agentes y búsqueda). & 2–3 h\\
Sáb 26 & Parte A: A4–A6. Parte B: B8–B19 (CSP, búsqueda local, juegos, utilidad). & 4 h\\
Dom 27 & Parte A: A7–A8. Parte B: B20–B32 (lógica proposicional y de primer orden). & 4 h\\
Lun 28 & Parte C: C1–C10 (formulación, heurísticas, A*, bidireccional, CSP, búsqueda local). & 4 h\\
Mar 29 & \textbf{E1 cronometrado} (examen 2025, 2 h) y corrección con las respuestas. Repaso de B17, C13 (expectimax) y C3 (A* inconsistente). & 3 h\\
Mié 30 & Mañana: E2–E5 y C11–C20 priorizando juegos y lógica (C11, C14, C17–C19). Tarde: \textbf{Parte D cronometrada} (2 h) y corrección. Noche: completar a mano la hoja A4 con lo que falló. & 5–6 h\\
Jue 1 & Examen. Antes: releer la hoja A4 y las plantillas de modelado (30 min). & —\\
\bottomrule
\end{tabularx}
\end{center}

\begin{nota}
\textbf{Estrategia.} En cada ejercicio escribe primero las definiciones que vas a usar
(por ejemplo, «$h$ es consistente si $h(n)\le c(n,a,n')+h(n')$»). En las tareas se evaluó
el procedimiento; en el examen también. Una demostración correcta con definiciones explícitas
vale más que un resultado sin justificación.
\end{nota}

\newpage
% ======================= FORMULARIO =======================
\section{Formulario de repaso}

\begin{formulario}
\textbf{Agentes.} Agente racional: para cada secuencia de percepciones elige la acción que
maximiza el desempeño \emph{esperado}, dada la evidencia y el conocimiento previo.
PEAS = \emph{Performance, Environment, Actuators, Sensors}. Propiedades del entorno:
observable (total/parcial), uno/varios agentes, determinista/estocástico, episódico/secuencial,
estático/dinámico, discreto/continuo, conocido/desconocido.
\end{formulario}

\begin{formulario}
\textbf{Problema de búsqueda:} estado inicial, $\textsc{Acciones}(s)$, $\textsc{Resultado}(s,a)$,
prueba de meta, costo $c(s,a,s')$. $b$ = factor de ramificación, $d$ = profundidad de la solución
más somera, $m$ = profundidad máxima, $C^*$ = costo óptimo, $\epsilon$ = costo mínimo de acción.

\smallskip
\centering
\renewcommand{\arraystretch}{1.2}
\begin{tabular}{lcccc}
\toprule
 & BFS & UCS & DFS & IDS \\
\midrule
Frontera & cola FIFO & cola de prioridad ($g$) & pila LIFO & pila con límite\\
¿Completa? & Sí$^a$ & Sí$^{a,b}$ & No (sí en finito, grafo) & Sí$^a$\\
¿Óptima? & Sí$^c$ & Sí & No & Sí$^c$\\
Tiempo & $O(b^d)$ & $O(b^{1+\lfloor C^*/\epsilon\rfloor})$ & $O(b^m)$ & $O(b^d)$\\
Espacio & $O(b^d)$ & $O(b^{1+\lfloor C^*/\epsilon\rfloor})$ & $O(bm)$ & $O(bd)$\\
\bottomrule
\end{tabular}

\raggedright\footnotesize $^a$ si $b$ es finito; $^b$ si todo costo $\ge\epsilon>0$; $^c$ si todos los costos son iguales.
Nodos generados por IDS: $(d)b+(d-1)b^2+\dots+(1)b^d$. Bidireccional: $O(b^{d/2})$.
\end{formulario}

\begin{formulario}
\textbf{Búsqueda informada.} Voraz: $f(n)=h(n)$. A*: $f(n)=g(n)+h(n)$. A* ponderada: $f=g+w\,h$.
\begin{itemize}[leftmargin=1.2em]
\item Admisible: $0\le h(n)\le h^*(n)$. \quad Consistente: $h(n)\le c(n,a,n')+h(n')$ y $h(\text{meta})=0$.
\item Consistente $\Rightarrow$ admisible. Con $h$ consistente, $f$ no decrece a lo largo de un camino.
\item A* en árbol + $h$ admisible $\Rightarrow$ óptimo. A* en grafo + $h$ consistente $\Rightarrow$ óptimo y sin reexpansiones.
\item Dominancia: $h_2\ge h_1$ (ambas admisibles) $\Rightarrow$ A* con $h_2$ no expande más nodos (salvo empates).
A* expande seguramente todo nodo con $f(n)<C^*$.
\item $\max(h_1,\dots,h_k)$ de admisibles es admisible (y consistente si todas lo son).
\item Factor de ramificación efectivo: $N+1=1+b^*+(b^*)^2+\dots+(b^*)^d$.
\end{itemize}
\end{formulario}

\begin{formulario}
\textbf{CSP} $=(X,D,C)$. Backtracking = DFS asignando una variable por nivel y verificando
restricciones. Revisión hacia adelante: al asignar $X$, borrar de los vecinos los valores
inconsistentes. Arco $X\to Y$ consistente: $\forall x\in D_X\ \exists y\in D_Y$ compatible.
AC-3: si $\textsc{Revise}(X,Y)$ borra algo, reencolar $(Z,X)$ para todo vecino $Z\ne Y$;
costo $O(cd^3)$ con $c$ arcos. MRV: variable con menos valores legales (desempate: mayor grado).
LCV: valor que elimina menos opciones a los vecinos. CSP con estructura de árbol: $O(nd^2)$.
\end{formulario}

\begin{formulario}
\textbf{Búsqueda local.} Hill climbing: moverse al mejor vecino; parar si ninguno mejora
(máximos locales, mesetas, crestas). Tabú: mejor vecino \emph{no tabú} aunque empeore.
Recocido simulado (maximizando): si $\Delta E=f(n)-f(s)>0$ aceptar; si no, aceptar con
probabilidad $e^{\Delta E/T}$; $T\leftarrow\alpha T$. Genético: población, aptitud, selección
(proporcional o torneo), cruza, mutación ($P_m\approx 1/L$). No free lunch: promediado sobre
todos los problemas, ningún algoritmo supera a otro.
\end{formulario}

\begin{formulario}
\textbf{Juegos.} $\textsc{Minimax}(s)=\textsc{Utilidad}(s)$ si es terminal; $\max_a$ en MAX; $\min_a$
en MIN. Tiempo $O(b^m)$, espacio $O(bm)$. Alfa-beta: $\alpha$ = mejor valor garantizado para MAX
en el camino actual, $\beta$ = mejor (menor) para MIN; con orden perfecto $O(b^{m/2})$, orden
aleatorio $\approx O(b^{3m/4})$. Expectimax: nodos de azar $=\sum_i p_i\,V(s_i)$. $n$ jugadores:
vectores de utilidad, cada jugador maximiza su componente. Evaluación lineal:
$\textsc{Eval}(s)=\sum_i w_i f_i(s)$.

\textbf{Utilidad.} Lotería $L=[p,A;\,1-p,B]$. Axiomas: ordenabilidad, transitividad,
continuidad, sustituibilidad, monotonicidad (y descomponibilidad en AIMA).
Teorema: existe $U$ con $U(A)\ge U(B)\iff A\succeq B$ y $U(L)=\sum_i p_iU(S_i)$.
\end{formulario}

\begin{formulario}
\textbf{Lógica proposicional.} $\alpha\ent\beta$ sii $M(\alpha)\subseteq M(\beta)$
sii $\alpha\imp\beta$ es válida sii $\alpha\wedge\neg\beta$ es insatisfacible.
Equivalencias: $\alpha\imp\beta\equiv\neg\alpha\vee\beta$;
$\alpha\sii\beta\equiv(\alpha\imp\beta)\wedge(\beta\imp\alpha)$;
De Morgan; contrapositiva $\alpha\imp\beta\equiv\neg\beta\imp\neg\alpha$;
distributividad $\alpha\vee(\beta\wedge\gamma)\equiv(\alpha\vee\beta)\wedge(\alpha\vee\gamma)$.

CNF: (1) eliminar $\sii$, (2) eliminar $\imp$, (3) meter $\neg$ (De Morgan, doble negación),
(4) distribuir $\vee$ sobre $\wedge$.
Resolución: $\dfrac{\ell_1\vee\dots\vee\ell_k\vee p,\quad \neg p\vee m_1\vee\dots\vee m_n}{\ell_1\vee\dots\vee\ell_k\vee m_1\vee\dots\vee m_n}$
(correcta y completa por refutación).
Horn: a lo más un literal positivo; definida: exactamente uno. Encadenamiento
hacia adelante/atrás: lineal en el tamaño de la KB; modus ponens es completo para Horn.
DPLL = backtracking + terminación temprana + símbolos puros + cláusulas unitarias.
\end{formulario}

\begin{formulario}
\textbf{Lógica de primer orden.} $\forall x\,(P(x)\imp Q(x))$ \;/\; $\exists x\,(P(x)\wedge Q(x))$.
$\neg\forall x\,P\equiv\exists x\,\neg P$; $\neg\exists x\,P\equiv\forall x\,\neg P$.
CNF: (1) eliminar $\sii,\imp$; (2) meter $\neg$; (3) estandarizar variables; (4) skolemizar
($\exists y$ dentro de $\forall x_1\dots x_k$ $\Rightarrow$ $F(x_1,\dots,x_k)$; sin $\forall$ externos
$\Rightarrow$ constante); (5) quitar $\forall$; (6) distribuir $\vee$ sobre $\wedge$.
Unificación: $\textsc{Unify}(p,q)=\theta$ tal que $p\theta=q\theta$; MGU = el más general;
verificación de ocurrencia: $x$ no unifica con un término que contenga a $x$.
Resolución de primer orden: resolver $p$ y $\neg q$ con $\theta=\textsc{Unify}(p,q)$ y aplicar
$\theta$ al resolvente. Prolog: cláusulas de Horn, resolución SLD en profundidad, de izquierda a
derecha y de arriba hacia abajo.
\end{formulario}

\newpage
% ======================= PARTE A =======================
\section{Parte A — Conceptos (Nivel 1)}

\begin{ejercicio}{A1 · Agentes e introducción \dif{1}}
\begin{enumerate}[label=\textbf{\arabic*.}]
\item Define \emph{agente racional}. ¿Racionalidad es lo mismo que omnisciencia? Da un ejemplo que
distinga ambos conceptos.
\item Enuncia los cuatro enfoques de la IA (pensar/actuar $\times$ humanamente/racionalmente).
¿Cuál adopta el curso y por qué es el más conveniente en ingeniería?
\item ¿Qué evalúa la prueba de Turing? ¿Qué pretende mostrar el argumento de la habitación
china de Searle? Distingue IA fuerte de IA débil.
\item Verdadero o falso, justificando en una línea \fuente{AIMA, ejercicios cap.\ 2}:
\begin{enumerate}[label=(\alph*)]
\item Un agente que solo percibe información parcial del estado no puede ser perfectamente racional.
\item Existen entornos de tarea en los que ningún agente puramente reflejo se comporta racionalmente.
\item Existe un entorno de tarea en el que todo agente es racional.
\item La entrada del programa de agente es la misma que la entrada de la función de agente.
\item Toda función de agente puede implementarse con alguna combinación programa/máquina.
\item Un agente que elige sus acciones uniformemente al azar puede ser racional en algún entorno determinista.
\item Un agente de póker perfectamente racional nunca pierde.
\end{enumerate}
\item ¿En qué se distingue un agente reflejo de un agente que planea? ¿Puede un agente reflejo ser racional?
\end{enumerate}
\end{ejercicio}

\begin{ejercicio}{A2 · Búsqueda no informada \dif{1}}
\begin{enumerate}[label=\textbf{\arabic*.}]
\item Enumera los componentes de un problema de búsqueda. ¿Qué es una solución? ¿Y una solución óptima?
\item Distingue \emph{estado} de \emph{nodo}, y \emph{grafo del espacio de estados} de
\emph{árbol de búsqueda}. ¿Por qué un grafo finito puede producir un árbol infinito?
\item ¿Qué cambia entre DFS, BFS, UCS y A*? Responde en términos de la frontera.
\item Verdadero o falso \fuente{AIMA, ejercicios cap.\ 3}:
(a) BFS es un caso particular de UCS. (b) DFS es un caso particular de búsqueda primero-el-mejor.
(c) UCS es un caso particular de A*. (d) BFS es completa aunque existan costos de acción iguales a cero.
(e) DFS siempre expande al menos tantos nodos como A* con heurística admisible.
\item ¿Por qué en UCS la prueba de meta debe hacerse al \emph{extraer} el nodo y no al \emph{generarlo}?
\item ¿Por qué IDS no desperdicia mucho tiempo a pesar de regenerar los niveles superiores?
\item Explica de dónde sale la cota $O(b^{1+\lfloor C^*/\epsilon\rfloor})$ de UCS. ¿Qué ocurre si $\epsilon\to 0$?
\end{enumerate}
\end{ejercicio}

\begin{ejercicio}{A3 · Heurísticas \dif{1}}
\begin{enumerate}[label=\textbf{\arabic*.}]
\item Define heurística admisible y consistente. ¿Cuál implica a cuál? Da un ejemplo de admisible no consistente.
\item ¿Qué es un problema relajado? ¿De qué relajaciones del 8-puzzle salen $h_1$ (fichas fuera de
lugar) y $h_2$ (Manhattan)?
\item ¿Qué significa que $h_2$ domine a $h_1$? ¿Por qué conviene usar la dominante?
\item ¿Es óptima la búsqueda voraz? ¿Es completa en búsqueda en árbol? ¿Y en grafo finito?
\item ¿Qué algoritmo es A* con $h\equiv 0$? ¿Qué hace A* con $h=h^*$?
\item Define el factor de ramificación efectivo $b^*$. ¿Qué valor ideal tendría?
\end{enumerate}
\end{ejercicio}

\begin{ejercicio}{A4 · Problemas de satisfacción de restricciones \dif{1}}
\begin{enumerate}[label=\textbf{\arabic*.}]
\item Define CSP. ¿Qué es una asignación consistente, completa y una solución?
\item Da un ejemplo de restricción unaria, binaria, global (de alto orden) y suave.
\item Diferencia entre revisión hacia adelante (\emph{forward checking}) y consistencia de arcos.
¿Cuál detecta más fallos? ¿Cuál cuesta más?
\item MRV, heurística de grado y LCV: ¿qué decide cada una? ¿Por qué para variables se elige
la que «falla primero» y para valores el que «falla al último»?
\item Sin fijar el orden de variables, el árbol de búsqueda ingenuo tiene $n!\,d^n$ hojas. Explica
por qué y por qué la conmutatividad permite reducirlo a $d^n$.
\item Después de ejecutar AC-3, ¿qué tres situaciones pueden presentarse?
\end{enumerate}
\end{ejercicio}

\begin{ejercicio}{A5 · Búsqueda local y metaheurísticas \dif{1}}
\begin{enumerate}[label=\textbf{\arabic*.}]
\item ¿En qué se distingue la búsqueda local de la búsqueda sistemática? ¿Por qué es incompleta?
\item Describe los tres obstáculos de hill climbing: máximos locales, mesetas y crestas.
\item Recocido simulado: ¿cómo se comporta con temperatura muy alta? ¿Y muy baja?
\item Búsqueda tabú: ¿para qué sirve la lista tabú? ¿Qué pasa si es muy corta o muy larga?
\item Enumera los elementos de un algoritmo genético. ¿Qué ocurre si $P_m$ es muy alta? ¿Y si es nula?
\item Enuncia informalmente el teorema \emph{No free lunch}. ¿Qué consecuencia práctica tiene?
\end{enumerate}
\end{ejercicio}

\begin{ejercicio}{A6 · Búsqueda con adversarios y utilidad \dif{1}}
\begin{enumerate}[label=\textbf{\arabic*.}]
\item Define formalmente un juego. ¿Qué es un juego de suma cero? ¿Y de información perfecta?
\item Complejidad en tiempo y espacio de minimax. ¿Por qué no se puede aplicar tal cual al ajedrez?
\item ¿Qué representan $\alpha$ y $\beta$? ¿Por qué la poda no cambia el valor de la raíz pero sí
puede dejar valores intermedios «incorrectos»?
\item ¿Cómo influye el orden de los movimientos en alfa-beta? Da la complejidad con orden perfecto.
\item ¿Qué propiedades debe tener una buena función de evaluación? ¿Qué es el efecto horizonte?
\item ¿Cuándo usar expectimax en vez de minimax? ¿Por qué en expectimax importa la
\emph{magnitud} de las utilidades y en minimax solo su \emph{orden}?
\item Enuncia los axiomas de racionalidad sobre preferencias y el teorema de utilidad esperada.
\end{enumerate}
\end{ejercicio}

\begin{ejercicio}{A7 · Lógica proposicional \dif{1}}
\begin{enumerate}[label=\textbf{\arabic*.}]
\item Distingue sintaxis, semántica y teoría de pruebas.
\item Define modelo, $KB\ent\alpha$, validez, satisfacibilidad y equivalencia lógica.
\item Justifica: $\alpha\ent\beta$ sii $(\alpha\imp\beta)$ es válida sii $(\alpha\wedge\neg\beta)$ es insatisfacible.
\item $\textsc{Tell}(f)$ puede responder «ya lo sabía», «no lo creo» o «aprendí algo». Expresa cada
respuesta en términos de $M(KB)$ y $M(f)$. Haz lo mismo para $\textsc{Ask}(f)$.
\item ¿Qué significa que un procedimiento de inferencia sea correcto (coherente) y completo?
¿Modus ponens es completo en general? ¿Y con cláusulas de Horn?
\item Define cláusula de Horn, cláusula definida y cláusula meta. ¿Es $A\vee B$ de Horn?
\item ¿Por qué la resolución se usa por refutación?
\end{enumerate}
\end{ejercicio}

\begin{ejercicio}{A8 · Lógica de primer orden \dif{1}}
\begin{enumerate}[label=\textbf{\arabic*.}]
\item ¿Qué agrega la LPO a la lógica proposicional? Enumera los tipos de términos y fórmulas.
\item ¿Por qué con $\forall$ se usa normalmente $\imp$ y con $\exists$ se usa $\wedge$? Muestra qué sale
mal con $\forall x\,(\p{Est}(x)\wedge\p{Listo}(x))$ y con $\exists x\,(\p{Est}(x)\imp\p{Listo}(x))$.
\item ¿Qué es la skolemización? ¿Por qué la función de Skolem depende de las variables
universales que \emph{rodean} al existencial?
\item ¿Qué es un unificador más general? ¿Para qué sirve la verificación de ocurrencia?
\item ¿Qué forma restringida de resolución usa Prolog? ¿Por qué puede ciclar aun cuando la
consulta sea consecuencia lógica del programa?
\end{enumerate}
\end{ejercicio}

\newpage
% ======================= PARTE B =======================
\section{Parte B — Procedimientos (Nivel 2)}

\subsection{Agentes y formulación}

\begin{ejercicio}{B1 · PEAS y propiedades del entorno \dif{2}\fuente{AIMA, ejercicios cap.\ 2}}
Para cada actividad, da la descripción PEAS y clasifica el entorno (observable, agentes,
determinista, episódico, estático, discreto):
(a) jugar fútbol; (b) comprar libros usados de IA por internet; (c) jugar un partido de tenis;
(d) practicar tenis contra una pared; (e) diagnóstico médico; (f) pujar por un artículo en una subasta.
\end{ejercicio}

\begin{ejercicio}{B2 · Misioneros y caníbales \dif{2}\fuente{AIMA, ejercicios cap.\ 3}}
Tres misioneros y tres caníbales están en una orilla de un río con un bote para una o dos
personas. Nunca puede haber en una orilla más caníbales que misioneros si hay al menos un misionero ahí.
\begin{enumerate}[label=(\alph*)]
\item Formula el problema: representación del estado, estado inicial, acciones, modelo de transición,
prueba de meta y costo.
\item ¿Cuántos estados tiene tu representación y cuántos son legales?
\item Encuentra una solución óptima (número de cruces). ¿Qué algoritmo garantiza la óptima?
\item ¿Por qué la gente encuentra difícil este problema aunque el espacio sea pequeño?
\end{enumerate}
\end{ejercicio}

\begin{ejercicio}{B3 · Tamaños de espacios de estados \dif{2}\fuente{AIMA §3.2, §3.6}}
\begin{enumerate}[label=(\alph*)]
\item ¿Cuántos estados alcanzables tiene el 8-puzzle? ¿Y el 15-puzzle? Explica el factor $1/2$.
\item 8 reinas, formulación incremental ingenua (colocar una reina en cualquier casilla vacía): cuenta
las secuencias posibles. Compárala con la formulación «una reina por columna, en la columna más
a la izquierda sin atacar» (2\,057 estados).
\item Estima cuántas partidas de gato existen (cota superior simple).
\end{enumerate}
\end{ejercicio}

\subsection{Búsqueda no informada e informada}

\begin{ejercicio}{B4 · Trazas en la gráfica $G_1$ \dif{2}\fuente{Clase / Tarea 1}}
\begin{center}\figGUno\end{center}
Arcos dirigidos con su costo; los valores $h$ se usan solo en B5. Convenciones: búsqueda en grafo,
prueba de meta al extraer, empates alfabéticos. Para cada algoritmo da el \textbf{orden de expansión},
el \textbf{camino devuelto} y su \textbf{costo}:
(a) BFS; (b) DFS; (c) IDS (búsqueda en árbol limitada, sin repetir nodos del camino actual;
indica cada iteración); (d) UCS, mostrando la frontera con sus valores $g$ en cada paso.
\end{ejercicio}

\begin{ejercicio}{B5 · Voraz, A* y propiedades de $h$ en $G_1$ \dif{2}\fuente{Clase}}
En la misma gráfica $G_1$:
\begin{enumerate}[label=(\alph*)]
\item Calcula $h^*(n)$ para todo nodo.
\item ¿Es $h$ admisible? ¿Es consistente? Verifica arco por arco.
\item Ejecuta la búsqueda voraz y A*, con la frontera y los valores $f$ en cada paso. ¿Cuál es óptima?
\end{enumerate}
\end{ejercicio}

\begin{ejercicio}{B6 · Rumania \dif{2}\fuente{AIMA Fig.\ 3.1 y 3.16}}
Usa el mapa de Rumania y las distancias en línea recta a Bucarest ($h_{SLD}$):

{\small
\begin{tabular}{@{}lll@{}}
Arad–Zerind 75 & Arad–Sibiu 140 & Arad–Timisoara 118\\
Zerind–Oradea 71 & Oradea–Sibiu 151 & Timisoara–Lugoj 111\\
Lugoj–Mehadia 70 & Mehadia–Drobeta 75 & Drobeta–Craiova 120\\
Craiova–Rimnicu 146 & Craiova–Pitesti 138 & Sibiu–Fagaras 99\\
Sibiu–Rimnicu 80 & Rimnicu–Pitesti 97 & Fagaras–Bucarest 211\\
Pitesti–Bucarest 101 & Bucarest–Urziceni 85 & Bucarest–Giurgiu 90
\end{tabular}

\smallskip
$h_{SLD}$: Arad 366, Bucarest 0, Craiova 160, Drobeta 242, Fagaras 176, Lugoj 244, Mehadia 241,
Oradea 380, Pitesti 100, Rimnicu 193, Sibiu 253, Timisoara 329, Zerind 374.}
\begin{enumerate}[label=(\alph*)]
\item UCS de Sibiu a Bucarest: ¿por qué no se detiene cuando genera Bucarest por primera vez?
\item A* de Lugoj a Bucarest: da la secuencia de nodos extraídos con $g$, $h$ y $f$.
\item Búsqueda voraz de Timisoara a Bucarest: camino y costo. Compáralo con el óptimo (536).
\item \emph{(AIMA)} Con $h_{SLD}$ como heurística, ¿qué le pasa a la búsqueda voraz \emph{en árbol} para ir de Iasi a Fagaras?
\end{enumerate}
\end{ejercicio}

\begin{ejercicio}{B7 · 8-puzzle, IDS y factor efectivo \dif{2}\fuente{AIMA §3.4.4, §3.6}}
\begin{center}
\begin{tabular}{ccc}
\puzzle{7}{2}{4}{5}{}{6}{8}{3}{1} & \puzzle{1}{2}{3}{4}{5}{6}{}{7}{8} & \puzzle{}{1}{2}{3}{4}{5}{6}{7}{8}\\
Inicio 1 & Inicio 2 & Meta
\end{tabular}
\end{center}
\begin{enumerate}[label=(\alph*)]
\item Calcula $h_1$ y $h_2$ para ambos estados iniciales (la meta es la de la derecha). Para el
inicio 1 el óptimo es 26: ¿qué tan informadas son?
\item Para el inicio 2 calcula $h_1$, $h_2$ y el costo óptimo exacto.
\item Con $b=10$ y $d=5$, compara el número de nodos generados por BFS y por IDS.
\item A* encuentra una solución a profundidad 5 generando 52 nodos. Plantea la ecuación de $b^*$ y
verifica que $b^*\approx 1.92$.
\end{enumerate}
\end{ejercicio}

\subsection{CSP}

\begin{ejercicio}{B8 · Criptoaritmética \dif{2}\fuente{AIMA Fig.\ 5.2}}
Formula como CSP el problema $\mathrm{TWO}+\mathrm{TWO}=\mathrm{FOUR}$ (letras distintas, dígitos
distintos, sin ceros iniciales): variables (incluye acarreos), dominios y restricciones (incluye la
restricción global y las de columna). Luego encuentra una solución razonando con las restricciones.
\end{ejercicio}

\begin{ejercicio}{B9 · AC-3 en el mapa de Australia \dif{2}\fuente{AIMA, ejercicios cap.\ 5}}
\begin{center}\figAustralia\end{center}
Dominios $\{R,G,B\}$ y restricciones $\ne$ en cada arista. Con la asignación parcial
$\{WA=G,\ V=R\}$, usa AC-3 para mostrar que la inconsistencia se detecta sin buscar.
Escribe cada llamada a \textsc{Revise} que elimina valores. Además: ¿cuántas soluciones tiene
el problema completo con 3 colores?
\end{ejercicio}

\begin{ejercicio}{B10 · Revisión hacia adelante en 4 reinas \dif{2}\fuente{Clase}}
Variables $Q_1,\dots,Q_4$ (columna $i$), dominio $\{1,2,3,4\}$ (fila). Aplica backtracking con
revisión hacia adelante, orden de variables fijo $Q_1,Q_2,Q_3,Q_4$ y valores en orden creciente.
Muestra los dominios después de cada asignación y cada retroceso, hasta hallar la primera solución.
\end{ejercicio}

\begin{ejercicio}{B11 · Heurísticas de ordenamiento \dif{2}\fuente{AIMA §5.3.1}}
En el mapa de Australia (3 colores):
(a) sin asignaciones, ¿qué variable elige MRV? ¿Y MRV con desempate por grado?
(b) con $\{WA=R,\ NT=G\}$, ¿qué variable elige MRV y por qué?
(c) con $\{WA=R,\ NT=G\}$ y la variable $Q$, ¿qué valor recomienda LCV?
\end{ejercicio}

\subsection{Búsqueda local}

\begin{ejercicio}{B12 · Hill climbing en 4 reinas \dif{2}\fuente{AIMA §4.1.1}}
\begin{center}\figReinas{2}{1}{4}{4}\end{center}
Formulación de estado completo: una reina por columna; un movimiento cambia de fila a una reina
dentro de su columna (12 sucesores). $h$ = número de pares de reinas que se atacan.
(a) Calcula $h$ del estado $(2,1,4,4)$ (fila de la reina en cada columna).
(b) Calcula $h$ para los 12 sucesores y aplica un paso de hill climbing de máximo descenso.
(c) Continúa hasta que el algoritmo se detenga. ¿Llega a una solución?
\end{ejercicio}

\begin{ejercicio}{B13 · Paisaje unidimensional \dif{2}\fuente{Clase, quiz X–Y–Z}}
\begin{center}\figPaisaje\end{center}
Se maximiza $f$; los vecinos del estado $i$ son $i\pm1$.
(a) Aplica hill climbing desde 0, 4, 10 y 13. ¿Dónde termina cada uno? ¿Cuál es el óptimo global?
(b) Aplica búsqueda tabú desde el estado 2, con lista tabú de los 2 últimos estados visitados,
desempate hacia la derecha, 6 iteraciones. ¿Cuál es la mejor solución vista?
(c) ¿Con qué probabilidad aceptaría recocido simulado pasar de 7 a 8 con $T=10$, $T=1$, $T=0.1$?
\end{ejercicio}

\begin{ejercicio}{B14 · Algoritmo genético en 8 reinas \dif{2}\fuente{AIMA Fig.\ 4.6}}
Cada individuo es una cadena de 8 dígitos: el dígito $i$ es la fila (1 = abajo) de la reina en
la columna $i$. Aptitud = número de pares de reinas que \emph{no} se atacan (máximo 28).
\begin{enumerate}[label=(\alph*)]
\item Verifica que la aptitud de \texttt{24748552} es 24 y la de \texttt{32543213} es 11.
\item La población es \texttt{24748552} (24), \texttt{32752411} (23), \texttt{24415124} (20),
\texttt{32543213} (11). Calcula las probabilidades de selección proporcional a la aptitud.
\item Cruza \texttt{32752411} con \texttt{24748552} en el punto 3 y calcula la aptitud de los hijos.
\item Con torneo de tamaño 2, ¿cuál es la probabilidad de que el peor individuo sea padre en un torneo (muestreo con reemplazo)?
\item Con $P_m=1/8$ por gen, ¿qué probabilidad hay de que un hijo no sufra ninguna mutación?
\end{enumerate}
\end{ejercicio}

\subsection{Juegos y utilidad}

\begin{ejercicio}{B15 · Minimax y alfa-beta (árbol $T_1$) \dif{2}\fuente{Clase}}
\begin{center}\figTUno\end{center}
(a) Calcula el valor minimax de cada nodo y la jugada de MAX.
(b) Aplica alfa-beta de izquierda a derecha, indicando $(\alpha,\beta)$ al entrar a cada nodo.
¿Qué hojas no se evalúan?
\end{ejercicio}

\begin{ejercicio}{B16 · Alfa-beta y orden de movimientos \dif{2}\fuente{AIMA Fig.\ 6.2}}
\begin{center}\figTDos\end{center}
(a) Valor minimax y hojas podadas por alfa-beta de izquierda a derecha.
(b) Reordena las hojas (sin cambiar a qué nodo MIN pertenecen) para maximizar la poda.
¿Cuántas hojas se evalúan entonces? Relaciónalo con $O(b^{m/2})$.
\end{ejercicio}

\begin{ejercicio}{B17 · Expectiminimax \dif{2}\fuente{Clase}}
\begin{center}\figExp\end{center}
Nodos: MAX (triángulo hacia arriba), azar (círculo), MIN (triángulo hacia abajo).
(a) Calcula el valor de cada nodo y la jugada de MAX.
(b) ¿Qué jugaría MAX si tratara los nodos de azar como nodos MIN (minimax pesimista)?
(c) ¿Se podría podar algún subárbol en $c_3$ si se sabe que las hojas están en $[0,20]$?
\end{ejercicio}

\begin{ejercicio}{B18 · Juego de tres jugadores \dif{2}\fuente{AIMA Fig.\ 6.4, Clase}}
\begin{center}\figMaxN\end{center}
Cada hoja tiene el vector $(u_1,u_2,u_3)$. Juegan 1, luego 2, luego 3; cada uno maximiza su
componente. Propaga los vectores. ¿Qué obtiene cada jugador? ¿Por qué pueden surgir alianzas?
\end{ejercicio}

\begin{ejercicio}{B19 · Loterías y utilidad \dif{2}\fuente{AIMA §15.2–15.3, Clase}}
\begin{enumerate}[label=(\alph*)]
\item Con $U(x)=\sqrt{x}$ (en miles), compara $A=[0.8,\$4k;\,0.2,\$0]$ y $B=[1.0,\$3k]$. ¿Cuál es el
equivalente cierto de $A$?
\item Muestra que preferir $B\succ A$ y $C\succ D$ con $C=[0.2,\$4k;\,0.8,\$0]$ y $D=[0.25,\$3k;\,0.75,\$0]$
es incompatible con la maximización de utilidad esperada (paradoja de Allais).
\item Si $A\succ B\succ C$, ¿qué axioma garantiza que existe $p$ con $[p,A;1-p,C]\sim B$? Calcula $p$ si
$U(A)=10$, $U(B)=4$, $U(C)=0$.
\end{enumerate}
\end{ejercicio}

\subsection{Lógica proposicional}

\begin{ejercicio}{B20 · Validez y satisfacibilidad \dif{2}\fuente{AIMA, ejercicios cap.\ 7}}
Clasifica como válida, insatisfacible o contingente (usa tablas o equivalencias):
\begin{multicols}{2}
\begin{enumerate}[label=(\alph*)]
\item $(\p{Humo}\imp\p{Fuego})\imp(\neg\p{Fuego}\imp\neg\p{Humo})$
\item $(A\wedge\neg A)\imp B$
\item $(A\imp B)\wedge(B\imp A)\wedge\neg(A\sii B)$
\item $((A\vee B)\wedge(\neg A\vee C))\imp(B\vee C)$
\item $((A\imp B)\imp A)\imp A$
\item $(A\imp B)\imp(B\imp A)$
\item $(A\imp(B\imp C))\sii((A\wedge B)\imp C)$
\item $(P\imp Q)\imp(\neg P\imp\neg Q)$
\end{enumerate}
\end{multicols}
\end{ejercicio}

\begin{ejercicio}{B21 · ¿Vinculación correcta? \dif{2}\fuente{AIMA, ejercicios cap.\ 7}}
Decide si cada afirmación es correcta:
\begin{multicols}{2}
\begin{enumerate}[label=(\alph*)]
\item $\p{False}\ent\p{True}$
\item $\p{True}\ent\p{False}$
\item $(A\wedge B)\ent(A\sii B)$
\item $A\sii B\ent A\vee B$
\item $A\sii B\ent\neg A\vee B$
\item $(A\wedge B)\imp C\ent(A\imp C)\vee(B\imp C)$
\item $(C\vee(\neg A\wedge\neg B))\equiv((A\imp C)\wedge(B\imp C))$
\item $(A\vee B)\wedge(\neg C\vee\neg D\vee E)\ent(A\vee B)$
\item $(A\vee B)\wedge(\neg C\vee\neg D\vee E)\ent(A\vee B)\wedge(\neg D\vee E)$
\item $(A\vee B)\wedge\neg(A\imp B)$ es satisfacible
\item $(A\sii B)\wedge(\neg A\vee B)$ es satisfacible
\item $(A\sii B)\sii C$ tiene el mismo número de modelos que $A\sii B$
\end{enumerate}
\end{multicols}
\end{ejercicio}

\begin{ejercicio}{B22 · Conteo de modelos \dif{2}\fuente{AIMA, ejercicios cap.\ 7}}
Con el vocabulario $\{A,B,C,D\}$, ¿cuántos modelos tiene cada oración?
(a) $B\vee C$; (b) $\neg A\vee\neg B\vee\neg C\vee\neg D$; (c) $(A\imp B)\wedge A\wedge\neg B\wedge C\wedge D$;
(d) $(A\wedge B)\vee(B\wedge C)$.
\end{ejercicio}

\begin{ejercicio}{B23 · Mundo del Wumpus por revisión de modelos \dif{2}\fuente{AIMA §7.3, Fig.\ 7.5}}
El agente visitó $[1,1]$ (sin brisa) y $[2,1]$ (con brisa). Considera solo los símbolos
$P_{1,2},P_{2,2},P_{3,1}$ (pozos).
(a) Escribe la KB relevante. (b) Enumera los modelos de la KB. (c) Decide por revisión de modelos si
$KB\ent\neg P_{1,2}$ y si $KB\ent\neg P_{2,2}$.
\end{ejercicio}

\begin{ejercicio}{B24 · Conversión a CNF \dif{2}\fuente{AIMA §7.5.2, Clase}}
Convierte a CNF, mostrando cada paso:
(a) $B_{1,1}\sii(P_{1,2}\vee P_{2,1})$; (b) $(A\wedge B)\imp(C\sii D)$; (c) $\neg(P\imp(Q\wedge R))$;
(d) $\p{Verano}\imp(\p{Nieve}\imp\p{Raro})$; (e) $(\p{Verano}\imp\p{Nieve})\imp\p{Raro}$ (diapositiva). ¿Por qué (d) y (e) no son equivalentes?
\end{ejercicio}

\begin{ejercicio}{B25 · Resolución proposicional \dif{2}\fuente{AIMA §7.5, ejercicios cap.\ 7}}
\begin{enumerate}[label=(\alph*)]
\item KB $=\{B_{1,1}\sii(P_{1,2}\vee P_{2,1}),\ \neg B_{1,1}\}$. Demuestra $\neg P_{1,2}$ por resolución.
\item KB $=\{A\imp(B\vee C),\ A,\ \neg B,\ \neg C\}$ (diapositiva). Muestra que la KB es insatisfacible.
\item «Si llueve y no llevo paraguas, me mojo. Si me mojo, me enfermo. Llovió. No me enfermé.»
Demuestra por resolución que llevé paraguas.
\item \emph{El unicornio} (AIMA): «Si el unicornio es mítico, es inmortal; si no es mítico, es un
mamífero mortal. Si es inmortal o mamífero, tiene cuerno. Si tiene cuerno, es mágico.» ¿Se puede
demostrar que es mítico? ¿Mágico? ¿Que tiene cuerno? Justifica con resolución o con un modelo.
\end{enumerate}
\end{ejercicio}

\begin{ejercicio}{B26 · Encadenamiento hacia adelante y hacia atrás \dif{2}\fuente{AIMA Fig.\ 7.16}}
KB: $P\imp Q$;\; $L\wedge M\imp P$;\; $B\wedge L\imp M$;\; $A\wedge P\imp L$;\; $A\wedge B\imp L$;\; $A$;\; $B$.
(a) Ejecuta encadenamiento hacia adelante para la consulta $Q$ mostrando la agenda y los contadores
de premisas. (b) Ejecuta encadenamiento hacia atrás para $Q$. ¿Qué ciclo debe evitarse?
\end{ejercicio}

\begin{ejercicio}{B27 · DPLL \dif{2}\fuente{AIMA §7.6.1, Tarea 3}}
Aplica DPLL (orden de símbolos $A,B,C,D,E$; primero \textit{true}) a
\[(A\vee\neg B)\wedge(\neg A\vee C\vee D)\wedge(B\vee\neg C)\wedge(\neg D\vee\neg A)\wedge(E\vee A)\wedge(E\vee\neg C)\wedge(\neg A\vee\neg B\vee\neg C).\]
Indica en cada paso si usas símbolo puro, cláusula unitaria o ramificación. Da el modelo encontrado.
\end{ejercicio}

\subsection{Lógica de primer orden}

\begin{ejercicio}{B28 · Traducción a LPO \dif{2}\fuente{AIMA, ejercicios cap.\ 8}}
Vocabulario: $\p{Est}(x)$, $\p{Persona}(x)$, $\p{Toma}(x,c,s)$, $\p{Aprueba}(x,c,s)$, $\p{Ama}(x,y)$,
$\p{Compra}(x,y)$, $\p{Poliza}(y)$, $\p{Cara}(y)$, $\p{Afeita}(x,y)$, $\p{Barbero}(x)$, $\p{Hombre}(x)$,
$\p{Engaña}(x,y,t)$, $\p{Politico}(x)$, $\p{Tiempo}(t)$; constantes $\p{Griego}$, $\p{Frances}$, $\p{P2001}$.
\begin{enumerate}[label=(\alph*)]
\item No todos los estudiantes toman francés y griego en P2001.
\item Solo un estudiante reprobó griego en P2001.
\item Todo mundo ama a alguien. / Hay alguien a quien todo mundo ama. (¿Son equivalentes?)
\item Ninguna persona compra una póliza cara.
\item Hay un barbero que afeita a todos los hombres que no se afeitan a sí mismos.
\item Los políticos pueden engañar a algunas personas todo el tiempo y a todas las personas algún
tiempo, pero no a todas las personas todo el tiempo.
\item Nadie ama a quien no se ama a sí mismo.
\end{enumerate}
\end{ejercicio}

\begin{ejercicio}{B29 · CNF en primer orden \dif{2}\fuente{AIMA §9.5.1, Clase}}
Convierte a CNF (con skolemización):
\begin{enumerate}[label=(\alph*)]
\item $\forall x\,[\p{Persona}(x)\imp\exists y\,\p{Madre}(y,x)]$
\item $\exists x\,\forall y\,\p{Ama}(y,x)$
\item $\forall x\,[(\exists y\,P(x,y))\imp Q(x)]$
\item $\forall x\,[\forall y\,P(x,y)\imp\exists z\,Q(x,z)]$
\item $\neg\exists x\,[\p{Est}(x)\wedge\forall y\,(\p{Curso}(y)\imp\p{Aprueba}(x,y))]$
\item $\forall x\,[\forall y\,\p{Animal}(y)\imp\p{Ama}(x,y)]\imp[\exists y\,\p{Ama}(y,x)]$ (ejemplo de clase)
\end{enumerate}
\end{ejercicio}

\begin{ejercicio}{B30 · Unificación \dif{2}\fuente{AIMA §9.2.2, ejercicios cap.\ 9}}
Da el MGU o explica por qué no existe ($x,y,z,u$ variables; mayúsculas constantes o funciones):
\begin{multicols}{2}
\begin{enumerate}[label=(\alph*)]
\item $\p{Conoce}(\p{Juan},x)$, $\p{Conoce}(y,\p{Madre}(y))$
\item $\p{Conoce}(\p{Juan},x)$, $\p{Conoce}(x,\p{Isabel})$
\item $P(x,F(x))$, $P(F(y),y)$
\item $R(x,G(x),z)$, $R(F(y),G(F(A)),y)$
\item $\p{Ama}(x,\p{Padre}(x))$, $\p{Ama}(\p{Padre}(y),z)$
\item $P(A,x,F(G(y)))$, $P(z,F(z),F(u))$
\item $H(x,x)$, $H(F(y),G(z))$
\item $Q(F(x),G(y),x)$, $Q(F(G(z)),z,A)$
\end{enumerate}
\end{multicols}
En (b), ¿qué cambia si se estandarizan las variables por separado?
\end{ejercicio}

\begin{ejercicio}{B31 · Prolog \dif{2}\fuente{Tarea 3, Clase}}
\begin{lstlisting}
p(a). p(b). q(b). q(c).
r(X) :- p(X), q(X).
progenitor(juan, maria). progenitor(maria, luis). progenitor(luis, ana).
ancestro(X, Y) :- progenitor(X, Y).
ancestro(X, Y) :- progenitor(X, Z), ancestro(Z, Y).
\end{lstlisting}
(a) Traza la consulta \texttt{?- r(Z).} indicando los retrocesos. (b) ¿Qué respuestas da
\texttt{?- ancestro(juan, W).} y en qué orden? (c) ¿Qué pasa si la segunda regla de \texttt{ancestro}
se escribe \texttt{ancestro(X,Y) :- ancestro(X,Z), progenitor(Z,Y).} y se pide
\texttt{?- ancestro(ana, W).}?
\end{ejercicio}

\begin{ejercicio}{B32 · Encadenamiento en LPO: el crimen del coronel West \dif{2}\fuente{AIMA §9.3}}
«La ley dice que es un crimen que un estadounidense venda armas a naciones hostiles. El país Nono, un
enemigo de Estados Unidos, tiene algunos misiles, y todos sus misiles se los vendió el coronel West,
que es estadounidense.»
(a) Escribe la KB en cláusulas definidas de primer orden (misiles son armas; un enemigo de EE.\,UU.\ es hostil).
(b) Demuestra $\p{Criminal}(\p{West})$ por encadenamiento hacia adelante, indicando las sustituciones.
\end{ejercicio}

\newpage
% ======================= PARTE C =======================
\section{Parte C — Problemas tipo examen (Nivel 3)}

\subsection{Búsqueda y heurísticas}

\begin{ejercicio}{C1 · Formulación y cotas: robot recolector \dif{3}\fuente{estilo Tarea 1, Problema 1}}
Un robot está en una cuadrícula $n\times n$ sin obstáculos, en la casilla $H$ (base). Hay $k$ monedas
en casillas conocidas. El robot se mueve N, S, E, O con costo 1; al entrar a una casilla con moneda
la recoge automáticamente. Debe recoger todas las monedas y regresar a $H$.
\begin{enumerate}[label=(\alph*)]
\item Formula el espacio de estados (estado, estado inicial, meta, acciones, costo) y argumenta que
es completo: toda solución del problema original es un camino en tu espacio y viceversa.
\item Da el número de estados, una cota del factor de ramificación y una cota superior de la
profundidad de la solución óptima. Estima los nodos que genera BFS en árbol.
\item Compara BFS y DFS en memoria y optimalidad para este problema.
\item ¿Cuándo se repite un estado? Propón cómo evitar la exploración redundante sin perder la
optimalidad.
\item Considera $h_a(s)=\max_{c\in\text{restantes}}[d(p,c)+d(c,H)]$ (o $d(p,H)$ si no quedan
monedas) y $h_b(s)=$ número de monedas restantes, con $d$ la distancia Manhattan. Demuestra que
ambas son admisibles, que $h_a$ es consistente, y que ninguna domina a la otra. ¿Qué heurística usarías?
\end{enumerate}
\end{ejercicio}

\begin{ejercicio}{C2 · Consistencia, monotonía y no-reexpansión \dif{3}\fuente{Tarea 1, Problema 2}}
\begin{enumerate}[label=(\alph*)]
\item Demuestra que toda heurística consistente con $h(\text{meta})=0$ es admisible.
\item Demuestra que si $h$ es consistente, los valores de $f$ a lo largo de cualquier camino no decrecen.
\item Demuestra que con $h$ consistente, A* en grafo, al extraer un nodo $n$ por primera vez, ya
encontró un camino óptimo hacia $n$; concluye que nunca reexpande nodos y que es óptimo.
\item Da un ejemplo de heurística admisible e inconsistente (no puedes usar $G_2$).
\end{enumerate}
\end{ejercicio}

\begin{ejercicio}{C3 · Optimalidad de A* y el papel de la consistencia \dif{3}\fuente{Clase 4}}
\begin{center}\figGDos\end{center}
\begin{enumerate}[label=(\alph*)]
\item Demuestra que A* en \emph{árbol} con $h$ admisible es óptimo (argumento del ancestro en la frontera).
\item En $G_2$: ¿es $h$ admisible? ¿consistente? Ejecuta A* en \emph{grafo} sin reapertura de
nodos cerrados. ¿Qué costo devuelve? ¿Cuál es el óptimo?
\item Ejecuta A* en \emph{árbol} en $G_2$. ¿Qué costo devuelve? Explica la diferencia.
\item Propón dos correcciones distintas para que la búsqueda en grafo sea óptima.
\end{enumerate}
\end{ejercicio}

\begin{ejercicio}{C4 · Combinar heurísticas \dif{3}\fuente{AIMA §3.6.2, ejercicios cap.\ 3}}
Sean $h_1,h_2$ admisibles. Demuestra o refuta:
(a) $\max(h_1,h_2)$ es admisible; si ambas son consistentes, también lo es;
(b) $h_1+h_2$ es admisible;
(c) $\lambda h_1+(1-\lambda)h_2$ con $\lambda\in[0,1]$ es admisible;
(d) $2h_1$ es admisible;
(e) si $h_1$ es consistente, $\min(h_1,h_2)$ es consistente.
\end{ejercicio}

\begin{ejercicio}{C5 · Heurísticas del $N$-puzzle \dif{3}\fuente{Tarea 1, Problema 2}}
Para el 8-puzzle:
(a) Demuestra que $h_1$ (fichas fuera de lugar) y $h_2$ (Manhattan) son admisibles, identificando
el problema relajado cuyo costo óptimo es exactamente cada heurística.
(b) Demuestra que $h_2$ es consistente. ¿Lo es $h_1$?
(c) Demuestra que $h_2\ge h_1$ y explica por qué A* con $h_2$ expande a lo más los nodos que
expande con $h_1$ (salvo empates).
(d) Explica por qué la heurística de conflicto lineal $h_3=h_2+2\cdot lc$ sigue siendo admisible.
\end{ejercicio}

\begin{ejercicio}{C6 · A* ponderada \dif{3}\fuente{AIMA §3.5.5}}
Sea $h$ admisible y $w\ge1$. A* ponderada ordena la frontera por $f_w(n)=g(n)+w\,h(n)$.
Demuestra que, en búsqueda en árbol, el costo de la solución devuelta es a lo más $w\,C^*$.
¿Qué algoritmo se obtiene con $w=0$, $w=1$ y $w\to\infty$?
\end{ejercicio}

\begin{ejercicio}{C7 · Búsqueda bidireccional \dif{3}\fuente{Tarea 1, Problema 4}}
\begin{center}\figBidir\end{center}
Gráfica no dirigida. Se ejecuta Dijkstra bidireccional alternando expansiones (primero hacia
adelante desde $S$, luego hacia atrás desde $G$).
\begin{enumerate}[label=(\alph*)]
\item Ejecuta la versión ingenua que se detiene en cuanto un nodo ha sido \emph{expandido} en ambas
direcciones y devuelve el camino por ese nodo. ¿Qué costo devuelve? ¿Cuál es $C^*$?
\item Sea $\mu$ el costo de la mejor solución encontrada (actualizada cada vez que se genera un nodo
ya alcanzado desde el otro lado). Demuestra que si $\mu\le g^F_{\min}+g^B_{\min}$ entonces $\mu=C^*$.
Aplícalo a la gráfica.
\item ¿Por qué la condición $\mu\le f^F_{\min}+f^B_{\min}$ con $f=g+h$ \emph{no} es válida en general
con heurísticas admisibles?
\item Enuncia la prioridad de MM y explica por qué ninguna de las dos búsquedas expande nodos con
$g>C^*/2$.
\end{enumerate}
\end{ejercicio}

\begin{ejercicio}{C8 · Verdadero o falso con justificación \dif{3}\fuente{AIMA, ejercicios cap.\ 3–4}}
\begin{enumerate}[label=(\alph*)]
\item Si una torre de ajedrez puede moverse cualquier número de casillas en línea recta, la distancia
Manhattan es admisible para el problema de llevarla de $A$ a $B$ en el mínimo número de movimientos.
\item A* no sirve en robótica porque los estados y acciones son continuos.
\item Si $h$ es admisible, A* en grafo sin reapertura siempre es óptimo.
\item IDS con costos no uniformes es óptima.
\item Hill climbing con reinicios aleatorios es completo con probabilidad 1 (en un espacio finito).
\item Si en 8 reinas hill climbing tiene éxito con probabilidad $p=0.14$, el número esperado de reinicios es $\approx 7$.
\item La búsqueda en haz local con $k$ estados es equivalente a $k$ reinicios aleatorios en paralelo.
\end{enumerate}
\end{ejercicio}

\subsection{CSP y búsqueda local}

\begin{ejercicio}{C9 · Reyes en un tablero \dif{3}\fuente{variante de Tarea 2, Ejercicio 1}}
Se quieren colocar $k$ reyes en un tablero $n\times n$ de modo que ningún par se ataque (un rey
ataca las 8 casillas adyacentes).
\begin{enumerate}[label=(\alph*)]
\item Da dos formulaciones como CSP: (i) una variable por rey; (ii) una variable booleana por casilla.
En cada una: variables, dominios, restricciones (explícitas) y número de restricciones binarias.
\item ¿Qué formulación es preferible para backtracking con MRV? ¿Cómo rompes la simetría entre reyes en (i)?
\item Para colocar la \emph{mayor} cantidad de reyes, define un algoritmo de búsqueda local:
estado, $\textsc{Acciones}$, $\textsc{Resultado}$ y función objetivo. ¿Cuál es el máximo teórico?
\end{enumerate}
\end{ejercicio}

\begin{ejercicio}{C10 · Límites de la consistencia y estructura \dif{3}\fuente{AIMA §5.2, §5.5}}
\begin{enumerate}[label=(\alph*)]
\item Da un CSP que sea arco-consistente y no tenga solución.
\item Demuestra que AC-3 tarda $O(cd^3)$ con $c$ restricciones binarias y dominios de tamaño $d$.
\item Explica el algoritmo para CSP con estructura de árbol y demuestra que es $O(nd^2)$.
\item ¿Por qué un solucionador local que minimiza conflictos funciona muy bien en $n$-reinas y mal
en Sudoku difíciles?
\end{enumerate}
\end{ejercicio}

\subsection{Juegos}

\begin{ejercicio}{C11 · Minimax contra un oponente subóptimo \dif{3}\fuente{AIMA, ejercicios cap.\ 6}}
(a) Demuestra que, para cualquier árbol, la utilidad que obtiene MAX usando decisiones minimax contra
un MIN subóptimo nunca es menor que la que obtiene contra un MIN óptimo.
(b) Construye un árbol en el que MAX obtiene más que el valor minimax usando una estrategia
\emph{no} minimax contra un MIN subóptimo (y describe ese MIN).
\end{ejercicio}

\begin{ejercicio}{C12 · Corrección de la poda alfa-beta \dif{3}\fuente{Tarea 2, Ejercicio 4}}
Considera un camino $n_1,n_2,\dots,n_j$ desde la raíz, con $n_1$ un nodo MIN y $n_j$ un nodo MAX.
\begin{enumerate}[label=(\alph*)]
\item Escribe $n_1$ en función de $n_2$ y sus hermanos; luego $n_2$ en función de $n_3$, etc., y
obtén una expresión para $n_1$ que contenga a $n_j$.
\item Con $l_i$ (valor ya conocido de los hermanos a la izquierda de $n_i$) y $r_i$ (hermanos
a la derecha aún no explorados), reescribe la expresión.
\item Muestra que $n_j$ puede afectar a $n_1$ solo si $n_j<\min\{l_i : i \text{ impar}\}$ (nodos MIN
del camino). Concluye la regla de poda.
\item Repite para el caso en que $n_j$ es un nodo MIN.
\end{enumerate}
\end{ejercicio}

\begin{ejercicio}{C13 · Transformaciones de utilidad \dif{3}\fuente{AIMA, ejercicios cap.\ 6}}
\begin{enumerate}[label=(\alph*)]
\item Demuestra que si se aplica a todas las hojas una transformación estrictamente creciente $g$, la
jugada elegida por minimax no cambia.
\item Demuestra que expectimax es invariante ante transformaciones afines positivas $g(x)=ax+b$, $a>0$.
\item Da un árbol de expectimax en el que $g(x)=\sqrt{x}$ cambie la decisión.
\item Si las hojas están acotadas en $[L,U]$, explica cómo podar en un nodo de azar.
\end{enumerate}
\end{ejercicio}

\begin{ejercicio}{C14 · Función de evaluación en el gato \dif{3}\fuente{Tarea 2, Ejercicio 3}}
Usa $\textsc{Eval}(s)=3X_2+X_1-(3O_2+O_1)$, donde $X_n$ ($O_n$) es el número de líneas con exactamente
$n$ X (O) y ninguna O (X). Posición actual (juega X):
\begin{center}\gato{O}{}{}{}{X}{}{}{}{}\end{center}
(a) Calcula $\textsc{Eval}$ de la posición. (b) Usando la simetría, lista las jugadas distintas de X y
calcula su $\textsc{Eval}$ (profundidad 1). (c) Para cada jugada de X, calcula el mínimo sobre las respuestas
de O (profundidad 2) y elige la jugada minimax. (d) ¿Coinciden las decisiones a profundidad 1 y 2? ¿Por qué?
\end{ejercicio}

\subsection{Lógica}

\begin{ejercicio}{C15 · Vinculación, validez y corrección de la resolución \dif{3}\fuente{AIMA §7.4–7.5}}
\begin{enumerate}[label=(\alph*)]
\item Demuestra que $\alpha\ent\beta$ si y solo si $\alpha\wedge\neg\beta$ es insatisfacible.
\item Demuestra que la regla de resolución es correcta.
\item Demuestra que si $KB\ent\alpha$ entonces $KB\wedge\beta\ent\alpha$ para cualquier $\beta$ (monotonía).
\item Demuestra que el resolvente de dos cláusulas de Horn es una cláusula de Horn.
\end{enumerate}
\end{ejercicio}

\begin{ejercicio}{C16 · Clases de cláusulas y complejidad \dif{3}\fuente{Tarea 3, 3b}}
\begin{enumerate}[label=(\alph*)]
\item ¿Cuántas cláusulas de exactamente 3 literales semánticamente distintas hay con $n$ símbolos?
(Considera tautologías y literales repetidos.)
\item Explica por qué la resolución sobre 2-CNF termina en tiempo polinomial pero la cota análoga
no da un algoritmo polinomial para 3-CNF.
\item Explica por qué el encadenamiento hacia adelante con cláusulas de Horn es lineal en el tamaño de la KB.
\item ¿Se puede expresar $P\vee Q$ con cláusulas de Horn? ¿Qué se pierde en expresividad?
\end{enumerate}
\end{ejercicio}

\begin{ejercicio}{C17 · Rudolph y Scrooge \dif{3}\fuente{problema clásico de resolución}}
«Todo niño ama a Santa. Todo el que ama a Santa ama a cualquier reno. Rudolph es un reno y tiene
la nariz roja. Todo lo que tiene la nariz roja es raro o es un payaso. Ningún reno es un payaso.
Scrooge no ama nada que sea raro.» Demuestra por resolución que Scrooge no es un niño:
(a) traduce a LPO; (b) convierte a CNF agregando la negación de la conclusión; (c) da la refutación
indicando las sustituciones.
\end{ejercicio}

\begin{ejercicio}{C18 · ¿Mató la curiosidad al gato? \dif{3}\fuente{AIMA §9.5.3, Clase}}
A. Todo el que ama a todos los animales es amado por alguien. B. Todo el que mata a un animal no es
amado por nadie. C. Jack ama a todos los animales. D. Jack o Curiosidad mató al gato, que se llama Tuna.
Pregunta: ¿mató Curiosidad al gato? (Añade: los gatos son animales.)
Traduce, convierte a CNF y demuestra $\p{Mata}(\p{Curiosidad},\p{Tuna})$ por resolución.
\end{ejercicio}

\begin{ejercicio}{C19 · Cuantificadores y la cabeza de un animal \dif{3}\fuente{AIMA, ejercicios cap.\ 8–9}}
\begin{enumerate}[label=(\alph*)]
\item Demuestra que $\exists y\,\forall x\,\p{Ama}(x,y)\ent\forall x\,\exists y\,\p{Ama}(x,y)$ y
construye un modelo que muestre que la vinculación inversa no se cumple.
\item Variante de la Tarea 3: «Los perros son animales» implica «la cola de un perro es la cola de un
animal». Traduce con $\p{ColaDe}(t,x)$, $\p{Perro}(x)$, $\p{Animal}(x)$, niega la conclusión, pasa a CNF y
demuestra por resolución.
\end{enumerate}
\end{ejercicio}

\begin{ejercicio}{C20 · Preguntas de ensayo breve \dif{3}\fuente{Clase 1, AIMA cap.\ 1–2}}
(Responde en 5–8 líneas cada una.)
(a) ¿Por qué el curso adopta el enfoque de agente racional y no el de la prueba de Turing?
(b) Expón el argumento de la habitación china y la «respuesta del sistema».
(c) ¿Por qué el problema de la parada impone límites a lo que un agente puede decidir?
(d) Compara búsqueda sistemática, CSP y búsqueda local: ¿cuándo usarías cada una?
\end{ejercicio}

\newpage
% ======================= PARTE D =======================
\section{Parte D — Examen simulado (2 horas, 100 puntos)}

\begin{nota}
Resuélvelo sin notas y cronometrado. Distribución sugerida: 15 min por pregunta y 15 min de revisión.
Convenciones generales de la guía (empates alfabéticos, prueba de meta al extraer, poda con $\le$/$\ge$).
\end{nota}

\begin{ejercicio}{D1 · Formulación \ptst{12}}
Tienes una jarra de 4 litros y una de 3 litros, sin marcas, y una llave de agua. Quieres tener
exactamente 2 litros en alguna jarra.
(a) Formula el problema (estado, inicial, meta, acciones con precondiciones, costo). \pts{5}
(b) ¿Cuántos estados hay en tu representación y cuántos son alcanzables desde el inicial? \pts{3}
(c) Da una solución óptima en número de acciones y justifica su optimalidad. \pts{4}
\end{ejercicio}

\begin{ejercicio}{D2 · A* y consistencia \ptst{15}}
\begin{center}\figGTres\end{center}
(a) Calcula $h^*$ y decide si $h$ es admisible y si es consistente. \pts{4}
(b) Ejecuta A* en grafo sin reapertura, desempate alfabético. Da la tabla de extracciones y el costo. \pts{5}
(c) Repite con desempate que favorece a $B$. ¿Qué ocurre y por qué? \pts{3}
(d) Demuestra que si $h_1,h_2$ son consistentes, $\max(h_1,h_2)$ es consistente. \pts{3}
\end{ejercicio}

\begin{ejercicio}{D3 · CSP \ptst{13}}
Variables $A,B,C,D$ con dominio $\{1,2,3\}$ y restricciones $A<B$, $B<C$, $D\ne C$, $D>A$.
(a) Dibuja el grafo de restricciones. \pts{2}
(b) Ejecuta AC-3 y da los dominios finales. ¿Hace falta buscar? \pts{6}
(c) Sin AC-3, ¿qué variable elegiría MRV al inicio? ¿Y con desempate por grado? \pts{2}
(d) Explica la diferencia entre revisión hacia adelante y AC-3 con este ejemplo, tras asignar $A=1$. \pts{3}
\end{ejercicio}

\begin{ejercicio}{D4 · Alfa-beta \ptst{15}}
\begin{center}\figTTres\end{center}
(a) Calcula el valor minimax de la raíz. \pts{4}
(b) Aplica alfa-beta de izquierda a derecha. Indica $(\alpha,\beta)$ en cada nodo interno y las hojas no evaluadas. \pts{8}
(c) ¿Cuántas hojas se evaluarían con orden perfecto en un árbol uniforme con $b=2$ y $m=4$? \pts{3}
\end{ejercicio}

\begin{ejercicio}{D5 · Búsqueda local \ptst{10}}
(a) En recocido simulado (maximizando), el estado actual vale 10 y el vecino 7. Calcula la probabilidad
de aceptación con $T=5$ y $T=1$. ¿Qué pasa cuando $T\to 0$? \pts{4}
(b) En OneMax ($L=8$, aptitud = número de unos), cruza \texttt{11010010} y \texttt{01101101} en el punto
3, y da la aptitud de los hijos. Con $P_m=1/8$, ¿probabilidad de que un hijo no mute? \pts{4}
(c) ¿Por qué la búsqueda tabú puede escapar de un máximo local y hill climbing no? \pts{2}
\end{ejercicio}

\begin{ejercicio}{D6 · Lógica proposicional \ptst{15}}
(a) Demuestra que $((P\imp Q)\wedge(R\imp S)\wedge(P\vee R))\imp(Q\vee S)$ es válida. \pts{4}
(b) Demuestra por resolución que $\{P\imp Q,\ R\imp S,\ P\vee R\}\ent Q\vee S$. \pts{6}
(c) ¿Puede escribirse esta KB solo con cláusulas de Horn? ¿Qué implica para el uso de modus ponens? \pts{3}
(d) ¿Cuántos modelos tiene $P\vee R$ sobre $\{P,Q,R,S\}$? \pts{2}
\end{ejercicio}

\begin{ejercicio}{D7 · Lógica de primer orden \ptst{20}}
«Todo el que sabe leer es letrado. Los delfines no son letrados. Algunos delfines son inteligentes.»
Conclusión: «Algunos que son inteligentes no saben leer».
(a) Traduce premisas y conclusión a LPO. \pts{5}
(b) Niega la conclusión y convierte todo a CNF (indica la skolemización). \pts{5}
(c) Demuestra la conclusión por resolución, con sustituciones. \pts{6}
(d) Da el MGU de $\p{Conoce}(x,\p{Madre}(x))$ y $\p{Conoce}(\p{Padre}(y),z)$, y explica por qué
$P(x,x)$ y $P(F(y),y)$ no unifican. \pts{4}
\end{ejercicio}

\newpage
% ======================= PARTE E =======================
\section{Parte E — Estilo del examen real (2025)}

\begin{nota}
\textbf{Qué muestra el examen de 2025.} Cuatro preguntas de 25 puntos: (1) modelar un ambiente real y
reformularlo como juego competitivo; (2) A* en una gráfica pequeña con heurística admisible pero
inconsistente; (3) modelar un juego \emph{estocástico} y definir su operador de valor
(expectiminimax); (4) leer una KB en CNF de primer orden, derivar por resolución en 2–3 pasos y decidir
si otra conclusión se sigue. Se permitía \textbf{una hoja A4 de notas}. Predominan el modelado y la
explicación sobre las demostraciones largas: responde con estructura (listas, tablas) y justifica cada
decisión en una o dos líneas.
\end{nota}

\begin{ejercicio}{E1 · Examen 1 de 2025 (resuélvelo en 2 h) \dif{2}\fuente{Examen del profesor, 2025}}
\textbf{1. Agentes y ambientes competitivos (25).} Te contratan para diseñar un sistema de taxis
autónomos en la Ciudad de México; la tarea es decidir dónde colocar los taxis.
(a) \pts{15} Modela el ambiente: estados, perceptos, acciones y propiedades (observable/parcial,
determinista/estocástico, discreto/continuo). Sé breve.
(b) \pts{10} Existen compañías que compiten por el espacio. Reformula como juego competitivo y discute
qué algoritmo (minimax, expectimax) sería más adecuado.

\textbf{2. Búsqueda informada (25).} Aristas $S\to A\,(1)$, $S\to B\,(4)$, $A\to C\,(2)$, $B\to C\,(1)$, $C\to G\,(3)$;
$h(S)=5$, $h(A)=3$, $h(B)=2$, $h(C)=1$, $h(G)=0$.
(a) \pts{10} Ejecuta A* mostrando el orden de expansión. (b) \pts{10} ¿La solución es óptima? Justifica.
(c) \pts{5} La heurística es admisible pero no consistente. Explica por qué y qué cambios de
implementación preservan la optimalidad.

\textbf{3. Juegos con incertidumbre (25).} Connect-4 de $4\times4$: al dejar caer una ficha en una columna,
con probabilidad $\rho$ cae en esa columna y con probabilidad $1-\rho$ en una columna uniformemente aleatoria.
(a) \pts{10} Modela formalmente: estados, acciones, transición estocástica.
(b) \pts{15} Define un operador de valor que combine minimax y expectimax para este juego.

\textbf{4. Lógica: CNF y resolución (25).} KB:
1.\ $\neg\p{Estudiante}(x)\vee\neg\p{TomaIA}(x)\vee\p{Conoce}(x,\p{astar})$;\;
2.\ $\p{Algoritmo}(\p{astar})$;\; 3.\ $\p{Optimo}(\p{astar})$;\;
4.\ $\neg\p{Optimo}(a)\vee\p{UsaAdmisible}(a)$;\; 5.\ $\neg\p{UsaAdmisible}(a)\vee\p{Correcto}(a)$;\;
6.\ $\p{Estudiante}(\p{bob})$;\; 7.\ $\p{TomaIA}(\p{bob})$.
(a) \pts{10} Lectura en lenguaje natural de cada cláusula. (b) \pts{10} Por resolución, deriva
$\p{Conoce}(\p{bob},\p{astar})$ en 2–3 pasos. (c) \pts{5} ¿Se sigue $\p{Correcto}(\p{astar})$? Si sí, argumento
breve por resolución; si no, qué falta en la base.
\end{ejercicio}

\begin{ejercicio}{E2 · Modelar un ambiente real y su versión competitiva \dif{2}\fuente{estilo 2025, pregunta 1}}
Diseñas el agente que redistribuye bicicletas entre las estaciones de un sistema de bicis compartidas
(camiones que mueven bicis de estaciones llenas a vacías).
(a) Modela el ambiente: medida de desempeño, estados, perceptos, acciones y las seis propiedades
(observable, agentes, determinista, episódico, estático, discreto). ¿Es un problema de búsqueda, de
CSP o de búsqueda local? Justifica.
(b) Una empresa rival de scooters compite por los mismos usuarios y también decide dónde colocar sus
vehículos. Formúlalo como juego (jugadores, turnos, acciones, utilidades). ¿Es de suma cero? ¿Usarías
minimax, expectimax o expectiminimax? ¿Qué harías si el juego es demasiado grande para llegar a las hojas?
\end{ejercicio}

\begin{ejercicio}{E3 · A* con heurística inconsistente y reapertura \dif{2}\fuente{estilo 2025, pregunta 2}}
Aristas dirigidas: $S\to A\,(2)$, $S\to B\,(1)$, $A\to C\,(1)$, $B\to C\,(3)$, $C\to D\,(1)$, $D\to G\,(2)$.
Heurística: $h(S)=5$, $h(A)=4$, $h(B)=1$, $h(C)=1$, $h(D)=2$, $h(G)=0$.
(a) Calcula $h^*$; muestra que $h$ es admisible y encuentra los arcos que violan la consistencia.
(b) Ejecuta A* en grafo \emph{sin} reapertura. ¿Qué costo devuelve? ¿Es óptimo?
(c) Ejecuta A* \emph{con} reapertura de nodos cerrados. ¿Qué cambia?
(d) ¿Cuál es el mínimo número de valores de $h$ que hay que cambiar para que sea consistente (y siga admisible)? Propón uno.
\end{ejercicio}

\begin{ejercicio}{E4 · Gato con viento \dif{2}\fuente{estilo 2025, pregunta 3}}
En un gato estocástico, al marcar una casilla libre, con probabilidad $\rho$ la marca queda ahí y con
probabilidad $1-\rho$ queda en una casilla libre distinta, elegida uniformemente.
(a) Modela formalmente: estados, jugador en turno, acciones, modelo de transición $P(s'\mid s,a)$,
prueba terminal y utilidad.
(b) Escribe el operador expectiminimax para este juego con profundidad limitada y función de evaluación.
(c) En cierta posición MAX tiene dos jugadas. \emph{Atacar}: si la marca cae donde quiere, gana ($+1$);
si se desvía, pierde ($-1$). \emph{Bloquear}: si cae donde quiere, empata ($0$); si se desvía, cae con igual
probabilidad en una casilla que empata ($0$) o en una que pierde ($-1$). Calcula el valor de ambas jugadas
en función de $\rho$ y el umbral de $\rho$ a partir del cual conviene atacar. ¿Qué elegiría minimax (tratando
el azar como adversario)?
\end{ejercicio}

\begin{ejercicio}{E5 · KB en CNF: leer, derivar y detectar qué falta \dif{2}\fuente{estilo 2025, pregunta 4}}
KB: 1.\ $\neg\p{Robot}(x)\vee\neg\p{Bateria}(x)\vee\p{Mueve}(x)$;\; 2.\ $\neg\p{Mueve}(x)\vee\neg\p{EnRuta}(x)\vee\p{Entrega}(x)$;\;
3.\ $\p{Robot}(r_1)$;\; 4.\ $\p{Bateria}(r_1)$;\; 5.\ $\neg\p{Dron}(x)\vee\p{Vuela}(x)$;\; 6.\ $\neg\p{Vuela}(x)\vee\p{EvitaTrafico}(x)$;\;
7.\ $\p{EnRuta}(r_2)$.
(a) Lee cada cláusula en lenguaje natural (como implicación).
(b) Deriva $\p{Mueve}(r_1)$ por resolución en 2 pasos.
(c) ¿Se sigue $\p{Entrega}(r_1)$? ¿Y $\p{EvitaTrafico}(r_1)$? Si no, di qué falta y da un modelo de la KB donde sea falsa.
(d) Pasa a CNF: «Todo dron con batería baja regresa a la base» y «Todo pedido es entregado por algún
robot». ¿Cuál necesita skolemización?
\end{ejercicio}


\vfill
\begin{center}\textcolor{gris}{— Fin de la guía. ¡Éxito en el examen! —}\end{center}

\end{document}
